% ============================================================
\chapter{Rappels de Topologie G\'en\'erale et Langage Cat\'egoriel}
\label{ch:rappels}
% ============================================================

Qu'est-ce qu'un espace~? La question para\^it simple, presque na\"ive, et
pourtant elle a oblig\'e des g\'en\'erations de math\'ematiciens \`a repenser les
fondements de la g\'eom\'etrie. Au d\'ebut du XX\ieme{} si\`ecle, Felix Hausdorff,
Maurice Fr\'echet et d'autres ont forg\'e le concept d'\emph{espace topologique}
--- une structure si g\'en\'erale qu'elle englobe aussi bien les espaces
euclidiens familiers que des objets exotiques o\`u notre intuition vacille.

Ce chapitre est un point de d\'epart, mais pas un simple r\'echauff\'e d'un
cours de licence. Notre objectif est de revisiter les fondations avec un
\oe il neuf~: celui de la th\'eorie des cat\'egories. Nous allons reformuler les
constructions classiques --- continuit\'e, produits, quotients --- dans un
langage qui r\'ev\`ele leur structure profonde. Et surtout, nous allons
abandonner un pr\'ejug\'e tenace~: celui qui voudrait que tout espace
\og raisonnable\fg{} soit s\'epar\'e au sens de Hausdorff. Car les espaces non
s\'epar\'es sont au c\oe ur de ce cours, et c'est l\`a que les choses
deviennent vraiment int\'eressantes.

\begin{intuition}
Ce chapitre rappelle les notions fondamentales de topologie g\'en\'erale en les reformulant
dans le langage cat\'egoriel. L'objectif n'est pas de r\'ep\'eter un cours de L3, mais de
mettre en place le vocabulaire et les outils qui seront utilis\'es dans les chapitres
suivants, en insistant sur les aspects souvent n\'eglig\'es~: espaces non s\'epar\'es, r\^ole
de la cat\'egorie $\mathbf{Top}$ et ses propri\'et\'es, et les limites de l'intuition
hausdorffienne.
\end{intuition}

% ------------------------------------------------------------
\section{Espaces topologiques~: d\'efinitions et premiers exemples}
% ------------------------------------------------------------

Tout commence par trois axiomes. Trois r\`egles d'une simplicit\'e trompeuse
qui d\'eterminent quels sous-ensembles d'un ensemble $X$ m\'eritent d'\^etre
appel\'es \og ouverts\fg. C'est le choix de ces ouverts --- et non une notion
de distance --- qui d\'efinit la g\'eom\'etrie d'un espace. Hausdorff
lui-m\^eme, dans son \emph{Grundz\"uge der Mengenlehre} de 1914, a compris que
la topologie est avant tout une affaire de \emph{voisinages}, pas de
m\'etriques.

\begin{definition}[Espace topologique]
Un \emph{espace topologique} est un couple $(X,\tau)$ o\`u $X$ est un ensemble et
$\tau \subseteq \mathcal{P}(X)$ v\'erifie~:
\begin{enumerate}
  \item $\emptyset \in \tau$ et $X \in \tau$;
  \item $\tau$ est stable par r\'eunion quelconque~: si $(U_i)_{i\in I} \subseteq \tau$,
    alors $\bigcup_{i\in I} U_i \in \tau$;
  \item $\tau$ est stable par intersection finie~: si $U_1,\ldots,U_n \in \tau$, alors
    $U_1 \cap \cdots \cap U_n \in \tau$.
\end{enumerate}
Les \'el\'ements de $\tau$ sont appel\'es \emph{ouverts}. Les compl\'ementaires des ouverts
sont les \emph{ferm\'es}.
\end{definition}

L'asym\'etrie entre r\'eunions (quelconques) et intersections (finies) n'est
pas un d\'efaut de la d\'efinition~: elle refl\`ete une r\'ealit\'e profonde. Si
l'on autorisait les intersections infinies d'ouverts, la topologie
s'effondrerait sur elle-m\^eme~: dans $\R$, l'intersection de tous les
intervalles $(-1/n, 1/n)$ n'est que le singleton $\{0\}$, qui n'est pas
ouvert. Comprendre pourquoi cette asym\'etrie est n\'ecessaire, c'est d\'ej\`a
comprendre la topologie.

Pour se faire la main, consid\'erons les cas extr\^emes.

\begin{exemple}[Topologies extr\'emales]
Sur tout ensemble $X$~:
\begin{itemize}
  \item La \emph{topologie discr\`ete} $\tau = \mathcal{P}(X)$ rend tout sous-ensemble ouvert.
  \item La \emph{topologie grossi\`ere} (ou indiscr\`ete) $\tau = \{\emptyset, X\}$ ne poss\`ede
    que les ouverts triviaux.
\end{itemize}
\end{exemple}

Entre ces deux extr\^emes, il existe un espace minuscule mais
remarquablement riche, qui reviendra sans cesse dans ce cours.

\begin{exemple}[Espace de Sierpi\'nski]\label{ex:sierpinski}
L'espace de Sierpi\'nski est $\mathbb{S} = \{0,1\}$ muni de $\tau = \{\emptyset, \{1\},
\{0,1\}\}$. C'est un espace $T_0$ non $T_1$. Il joue un r\^ole fondamental~: le singleton
$\{0\}$ est ferm\'e mais non ouvert, et $\{1\}$ est ouvert mais non ferm\'e.
\end{exemple}

Cet espace de deux points est en r\'ealit\'e le \og classifieur des
ouverts\fg~: nous verrons dans les exercices que les ouverts d'un espace
$X$ sont en bijection avec les applications continues de $X$ vers
$\mathbb{S}$. C'est la premi\`ere manifestation d'un th\`eme r\'ecurrent~: les
petits espaces non s\'epar\'es encodent de l'information.

\begin{exemple}[Topologie cofinie]
Sur un ensemble infini $X$, la \emph{topologie cofinie} est
$\tau = \{\emptyset\} \cup \{U \subseteq X : X \setminus U \text{ est fini}\}$.
C'est un espace $T_1$ non $T_2$.
\end{exemple}

% ------------------------------------------------------------
\section{Applications continues et hom\'eomorphismes}
% ------------------------------------------------------------

En topologie, la notion de distance est absente, mais celle de
\og proximit\'e\fg{} persiste \`a travers les ouverts. La continuit\'e, qui en
analyse se formule avec des $\varepsilon$ et des $\delta$, prend ici une
forme plus \'epur\'ee~: une application est continue si elle respecte la
structure des ouverts, c'est-\`a-dire si l'image r\'eciproque de tout ouvert
est ouvert.

\begin{definition}[Application continue]
Soient $(X,\tau_X)$ et $(Y,\tau_Y)$ deux espaces topologiques. Une application
$f\colon X \to Y$ est \emph{continue} si pour tout $V \in \tau_Y$, on a
$f^{-1}(V) \in \tau_X$.
\end{definition}

Cette d\'efinition, d'une \'el\'egance laconique, admet plusieurs
reformulations \'equivalentes qui sont chacune utiles dans des contextes
diff\'erents.

\begin{proposition}[Caract\'erisations de la continuit\'e]
Soit $f\colon X \to Y$. Les conditions suivantes sont \'equivalentes~:
\begin{enumerate}
  \item $f$ est continue.
  \item Pour tout ferm\'e $F$ de $Y$, $f^{-1}(F)$ est ferm\'e dans $X$.
  \item Pour tout $A \subseteq X$, $f(\overline{A}) \subseteq \overline{f(A)}$.
  \item Pour tout $x \in X$ et tout voisinage $V$ de $f(x)$, $f^{-1}(V)$ est un
    voisinage de $x$.
\end{enumerate}
\end{proposition}

\begin{proof}
$(1) \Leftrightarrow (2)$~: imm\'ediat par passage au compl\'ementaire.

$(1) \Rightarrow (3)$~: Soit $A \subseteq X$ et $x \in \overline{A}$. Soit $V$ un ouvert
de $Y$ contenant $f(x)$. Alors $f^{-1}(V)$ est un ouvert contenant $x$, donc
$f^{-1}(V) \cap A \neq \emptyset$, d'o\`u $V \cap f(A) \neq \emptyset$. Ainsi
$f(x) \in \overline{f(A)}$.

$(3) \Rightarrow (2)$~: Si $F$ est ferm\'e dans $Y$ et $A = f^{-1}(F)$, alors
$f(\overline{A}) \subseteq \overline{f(A)} \subseteq \overline{F} = F$, donc
$\overline{A} \subseteq f^{-1}(F) = A$, ce qui montre que $A$ est ferm\'e.

$(1) \Leftrightarrow (4)$~: exercice.
\end{proof}

Quand deux espaces partagent la m\^eme \og forme topologique\fg, on dit qu'ils
sont hom\'eomorphes. C'est la notion d'\'equivalence en topologie~:
l'analogue de l'isom\'etrie en g\'eom\'etrie m\'etrique, ou de l'isomorphisme en
alg\`ebre.

\begin{definition}[Hom\'eomorphisme]
Une application $f\colon X \to Y$ est un \emph{hom\'eomorphisme} si $f$ est bijective et
si $f$ et $f^{-1}$ sont continues. On dit alors que $X$ et $Y$ sont \emph{hom\'eomorphes},
not\'e $X \cong Y$.
\end{definition}

% ------------------------------------------------------------
\section{Axiomes de s\'eparation}
% ------------------------------------------------------------

Voici l'un des th\`emes centraux de ce cours. Les axiomes de s\'eparation
classent les espaces topologiques selon leur capacit\'e \`a \og distinguer\fg{}
les points. Dans un cours classique, on se concentre sur les espaces de
Hausdorff ($T_2$) et on traite les axiomes inf\'erieurs comme des
curiosit\'es. Mais dans ce cours, les espaces $T_0$ non $T_1$ seront nos
compagnons quotidiens --- car ce sont pr\'ecis\'ement ceux qui apparaissent
en th\'eorie des domaines, en g\'eom\'etrie alg\'ebrique, et dans l'\'etude des
espaces spectraux.

\begin{definition}[Axiomes $T_0$, $T_1$, $T_2$]
Soit $(X,\tau)$ un espace topologique.
\begin{itemize}
  \item $X$ est \emph{$T_0$} (ou de Kolmogorov) si pour tous $x \neq y$ dans $X$, il
    existe un ouvert contenant l'un mais pas l'autre.
  \item $X$ est \emph{$T_1$} si pour tous $x \neq y$, il existe un ouvert contenant $x$
    mais pas $y$ (et par sym\'etrie, un ouvert contenant $y$ mais pas $x$).
  \item $X$ est \emph{$T_2$} (ou de Hausdorff) si pour tous $x \neq y$, il existe des
    ouverts disjoints $U \ni x$ et $V \ni y$.
\end{itemize}
\end{definition}

\begin{proposition}
On a les implications~: $T_2 \Rightarrow T_1 \Rightarrow T_0$, et aucune implication
inverse ne tient en g\'en\'eral.
\end{proposition}

La diff\'erence entre $T_0$ et $T_1$ est subtile mais fondamentale. Dans un
espace $T_0$, on peut \og distinguer\fg{} deux points, mais de mani\`ere
asym\'etrique~: l'un voit un ouvert que l'autre ne voit pas, mais pas
forc\'ement l'inverse. L'espace de Sierpi\'nski en est l'exemple type.

\begin{remarque}
Un espace est $T_1$ si et seulement si tout singleton est ferm\'e. Cette caract\'erisation
ne s'\'etend pas aux espaces $T_0$~: dans l'espace de Sierpi\'nski, $\{0\}$ est ferm\'e
mais $\{1\}$ ne l'est pas.
\end{remarque}

Au-del\`a de $T_2$, la hi\'erarchie se poursuit avec les axiomes de
r\'egularit\'e et de normalit\'e, qui portent sur la s\'eparation des ferm\'es.

\begin{definition}[Axiomes sup\'erieurs]
\begin{itemize}
  \item $X$ est \emph{r\'egulier} si pour tout ferm\'e $F$ et tout $x \notin F$, il existe
    des ouverts disjoints s\'eparant $x$ et $F$. $X$ est \emph{$T_3$} s'il est r\'egulier
    et~$T_1$.
  \item $X$ est \emph{normal} si pour tous ferm\'es disjoints $F_1, F_2$, il existe des
    ouverts disjoints les s\'eparant. $X$ est \emph{$T_4$} s'il est normal et~$T_1$.
\end{itemize}
\end{definition}

% ------------------------------------------------------------
\section{Constructions topologiques de base}
% ------------------------------------------------------------

Comment fabriquer de nouveaux espaces \`a partir d'espaces existants~? Trois
constructions fondamentales --- le produit, le quotient, et le sous-espace
--- permettent de b\^atir tout l'\'edifice de la topologie. Ce qui est
remarquable, c'est que chacune de ces constructions se caract\'erise par une
\emph{propri\'et\'e universelle}, un th\`eme que nous d\'evelopperons pleinement
dans la section cat\'egorielle.

\begin{definition}[Topologie produit]
Soit $(X_i,\tau_i)_{i\in I}$ une famille d'espaces topologiques. La \emph{topologie
produit} sur $\prod_{i\in I} X_i$ est la topologie la moins fine rendant continues
toutes les projections $\pi_j\colon \prod_{i\in I} X_i \to X_j$.
Une base de cette topologie est form\'ee des ensembles
$\prod_{i\in I} U_i$ o\`u $U_i \in \tau_i$ et $U_i = X_i$ sauf pour un nombre fini
d'indices.
\end{definition}

\begin{definition}[Topologie quotient]
Soit $f\colon X \to Y$ une surjection, o\`u $X$ est un espace topologique. La
\emph{topologie quotient} sur $Y$ est
$\tau_Y = \{V \subseteq Y : f^{-1}(V) \in \tau_X\}$.
C'est la topologie la plus fine rendant $f$ continue.
\end{definition}

\begin{definition}[Sous-espace]
Si $(X,\tau)$ est un espace topologique et $A \subseteq X$, la \emph{topologie induite}
(ou trace) sur $A$ est $\tau_A = \{U \cap A : U \in \tau\}$.
\end{definition}

% ------------------------------------------------------------
\section{Le langage cat\'egoriel}
% ------------------------------------------------------------

Nous arrivons maintenant \`a l'outil qui va transformer notre mani\`ere de
penser la topologie. La th\'eorie des cat\'egories, n\'ee en 1945 sous la plume
de Samuel Eilenberg et Saunders Mac Lane, est souvent d\'ecrite comme
\og les math\'ematiques des math\'ematiques\fg~: elle fournit un langage pour
parler des structures math\'ematiques et des relations entre elles, en
s'abstrayant de la nature des objets \'etudi\'es.

Pourquoi introduire ce formalisme dans un cours de topologie~? Parce que
les constructions topologiques que nous venons de voir --- produits,
quotients, sous-espaces --- sont des cas particuliers de \emph{limites} et
\emph{colimites} cat\'egorielles. En les formulant ainsi, on comprend
pourquoi ces constructions fonctionnent, et surtout pourquoi certaines
autres ne fonctionnent pas.

\begin{definition}[Cat\'egorie]
Une \emph{cat\'egorie} $\mathcal{C}$ consiste en~:
\begin{itemize}
  \item une classe $\mathrm{Ob}(\mathcal{C})$ d'\emph{objets};
  \item pour chaque paire d'objets $A,B$, un ensemble $\mathrm{Hom}_{\mathcal{C}}(A,B)$
    de \emph{morphismes};
  \item une loi de composition associative
    $\circ\colon \mathrm{Hom}(B,C) \times \mathrm{Hom}(A,B) \to \mathrm{Hom}(A,C)$;
  \item pour chaque objet $A$, un morphisme identit\'e $\mathrm{id}_A \in \mathrm{Hom}(A,A)$,
    neutre pour la composition.
\end{itemize}
\end{definition}

Voici les cat\'egories que nous utiliserons tout au long de ce cours.

\begin{exemple}[Cat\'egories topologiques fondamentales]
\begin{itemize}
  \item $\mathbf{Set}$~: ensembles et fonctions.
  \item $\mathbf{Top}$~: espaces topologiques et applications continues.
  \item $\mathbf{Top}_0$~: sous-cat\'egorie pleine de $\mathbf{Top}$ form\'ee des
    espaces~$T_0$.
  \item $\mathbf{Pos}$~: ensembles partiellement ordonn\'es et fonctions croissantes.
  \item $\mathbf{Lat}$~: treillis et morphismes de treillis.
\end{itemize}
\end{exemple}

Pour relier les cat\'egories entre elles, on utilise des \emph{foncteurs}~:
des \og traducteurs\fg{} qui pr\'eservent la structure cat\'egorielle.

\begin{definition}[Foncteur]
Un \emph{foncteur} $F\colon \mathcal{C} \to \mathcal{D}$ entre deux cat\'egories associe~:
\begin{itemize}
  \item \`a chaque objet $A$ de $\mathcal{C}$, un objet $F(A)$ de $\mathcal{D}$;
  \item \`a chaque morphisme $f\colon A \to B$, un morphisme $F(f)\colon F(A) \to F(B)$;
\end{itemize}
de sorte que $F(\mathrm{id}_A) = \mathrm{id}_{F(A)}$ et $F(g \circ f) = F(g) \circ F(f)$.
Un foncteur est \emph{contravariant} s'il renverse le sens des fl\`eches.
\end{definition}

Et pour comparer deux foncteurs entre eux, on utilise des
\emph{transformations naturelles} --- l'une des id\'ees les plus profondes
d'Eilenberg et Mac Lane, qui disaient d'ailleurs avoir invent\'e les
cat\'egories pour pouvoir d\'efinir les transformations naturelles.

\begin{definition}[Transformation naturelle]
Soient $F, G\colon \mathcal{C} \to \mathcal{D}$ deux foncteurs. Une \emph{transformation
naturelle} $\eta\colon F \Rightarrow G$ est une famille de morphismes
$(\eta_A\colon F(A) \to G(A))_{A \in \mathrm{Ob}(\mathcal{C})}$ telle que pour tout
morphisme $f\colon A \to B$ dans $\mathcal{C}$, le diagramme suivant commute~:
\[
\begin{tikzcd}
F(A) \arrow[r, "\eta_A"] \arrow[d, "F(f)"'] & G(A) \arrow[d, "G(f)"] \\
F(B) \arrow[r, "\eta_B"'] & G(B)
\end{tikzcd}
\]
\end{definition}

Le concept le plus important pour ce cours est celui d'\emph{adjonction},
qui capture la dualit\'e entre \og construire librement\fg{} et \og oublier la
structure\fg.

\begin{definition}[Adjonction]
Soient $F\colon \mathcal{C} \to \mathcal{D}$ et $G\colon \mathcal{D} \to \mathcal{C}$
deux foncteurs. On dit que $F$ est \emph{adjoint \`a gauche} de $G$ (not\'e $F \dashv G$)
s'il existe une bijection naturelle
\[
  \mathrm{Hom}_{\mathcal{D}}(F(A), B) \cong \mathrm{Hom}_{\mathcal{C}}(A, G(B))
\]
pour tous objets $A$ de $\mathcal{C}$ et $B$ de $\mathcal{D}$.
\end{definition}

Les adjonctions sont partout en topologie. Chaque fois que nous
rencontrerons une construction topologique, nous pourrons la comprendre
comme un adjoint \`a gauche ou \`a droite de quelque chose. Voici pourquoi
c'est si utile.

\begin{theoreme}[Propri\'et\'e universelle de l'adjonction]
Si $F \dashv G$, alors~:
\begin{enumerate}
  \item $F$ pr\'eserve les colimites.
  \item $G$ pr\'eserve les limites.
  \item Il existe des transformations naturelles $\eta\colon \mathrm{Id}_\mathcal{C}
    \Rightarrow GF$ (unit\'e) et $\varepsilon\colon FG \Rightarrow \mathrm{Id}_\mathcal{D}$
    (co-unit\'e) satisfaisant les identit\'es triangulaires.
\end{enumerate}
\end{theoreme}

\begin{proof}
C'est un r\'esultat classique de th\'eorie des cat\'egories. L'unit\'e $\eta_A$ correspond
\`a l'image de $\mathrm{id}_{F(A)}$ par la bijection
$\mathrm{Hom}_\mathcal{D}(F(A),F(A)) \cong \mathrm{Hom}_\mathcal{C}(A,GF(A))$.
La co-unit\'e $\varepsilon_B$ correspond \`a l'image de $\mathrm{id}_{G(B)}$. La
pr\'eservation des (co)limites d\'ecoule du fait qu'un adjoint \`a gauche (resp.\ droite)
commute aux colimites (resp.\ limites), car les Hom-foncteurs transforment les colimites
en limites.
\end{proof}

% ------------------------------------------------------------
\section{$\mathbf{Top}$ comme cat\'egorie~: propri\'et\'es fondamentales}
% ------------------------------------------------------------

Munis de ce langage, que pouvons-nous dire de la cat\'egorie $\mathbf{Top}$
elle-m\^eme~? La bonne nouvelle est qu'elle est tr\`es riche~: elle admet
toutes les limites et colimites imaginables. La mauvaise nouvelle --- et
c'est l\`a que les ennuis commencent --- est qu'elle n'est pas
cart\'esienne ferm\'ee.

\begin{theoreme}[$\mathbf{Top}$ est compl\`ete et cocompl\`ete]
La cat\'egorie $\mathbf{Top}$ admet toutes les petites limites et toutes les petites
colimites. En particulier~:
\begin{itemize}
  \item Les produits existent (topologie produit).
  \item Les \'egalisateurs existent (topologie induite).
  \item Les coproduits existent (somme topologique disjointe).
  \item Les co\'egalisateurs existent (topologie quotient).
\end{itemize}
\end{theoreme}

Le foncteur d'oubli $U\colon \mathbf{Top} \to \mathbf{Set}$ --- qui
\og oublie\fg{} la topologie et ne garde que l'ensemble sous-jacent --- est
remarquablement bien dot\'e~: il poss\`ede \`a la fois un adjoint \`a gauche et
un adjoint \`a droite, ce qui est relativement rare.

\begin{proposition}[Le foncteur d'oubli]
Le foncteur d'oubli $U\colon \mathbf{Top} \to \mathbf{Set}$ qui associe \`a un espace
son ensemble sous-jacent poss\`ede un adjoint \`a gauche (topologie discr\`ete) et un
adjoint \`a droite (topologie grossi\`ere). En notant $D$ et $I$ ces adjoints~:
\[
  D \dashv U \dashv I.
\]
\end{proposition}

\begin{proof}
Pour tout espace $Y$ et tout ensemble $S$, une fonction $f\colon S \to U(Y)$ est
automatiquement continue lorsqu'on munit $S$ de la topologie discr\`ete, car tout
sous-ensemble de $S$ est ouvert. Ainsi $\mathrm{Hom}_\mathbf{Top}(D(S),Y) \cong
\mathrm{Hom}_\mathbf{Set}(S,U(Y))$, d'o\`u $D \dashv U$.

De m\^eme, toute fonction $g\colon U(X) \to S$ est continue lorsqu'on munit $S$ de la
topologie grossi\`ere, car les seules pr\'eimages \`a v\'erifier sont $\emptyset$ et $S$.
D'o\`u $U \dashv I$.
\end{proof}

Mais voici le probl\`eme qui va motiver une grande partie de ce cours~:
$\mathbf{Top}$ \'echoue \`a \^etre cart\'esienne ferm\'ee.

\begin{theoreme}[$\mathbf{Top}$ est cart\'esienne ferm\'ee?]
La cat\'egorie $\mathbf{Top}$ n'est \emph{pas} cart\'esienne ferm\'ee. Plus pr\'ecis\'ement,
la topologie compacte-ouverte sur $\mathrm{Hom}(X,Y)$ ne donne pas en g\'en\'eral un
foncteur exponentiel adjoint au produit. Cependant, la restriction aux espaces
localement compacts de Hausdorff donne une sous-cat\'egorie cart\'esienne ferm\'ee.
\end{theoreme}

\begin{attention}
Le fait que $\mathbf{Top}$ ne soit pas cart\'esienne ferm\'ee est une limitation
fondamentale qui motive la recherche de cat\'egories topologiques plus agr\'eables,
comme les espaces de convergence (Chapitre~7) ou les espaces compactement engendr\'es.
\end{attention}

% ------------------------------------------------------------
\section{Le foncteur $\mathcal{O}$ et le treillis des ouverts}
% ------------------------------------------------------------

Changeons de point de vue. Au lieu d'\'etudier un espace \`a travers ses
\emph{points}, \'etudions-le \`a travers ses \emph{ouverts}. Cette id\'ee, qui
para\^it anodine, va nous mener vers la th\'eorie des locales
(Chapitre~3) et constitue l'un des ponts les plus profonds entre topologie
et logique.

L'ensemble des ouverts d'un espace topologique forme un treillis avec une
structure tr\`es particuli\`ere~: c'est un \emph{cadre} (\emph{frame}).

\begin{definition}[Treillis des ouverts]
Pour un espace topologique $X$, le \emph{treillis des ouverts} $\mathcal{O}(X)$ est
l'ensemble des ouverts de $X$, ordonn\'e par l'inclusion, muni des op\'erations~:
\begin{itemize}
  \item $\bigvee_{i\in I} U_i = \bigcup_{i\in I} U_i$ (sup quelconque);
  \item $U \wedge V = U \cap V$ (inf fini);
  \item $\bot = \emptyset$, $\top = X$.
\end{itemize}
C'est un \emph{cadre} (\emph{frame}), c'est-\`a-dire un treillis complet satisfaisant la
loi de distributivit\'e infinie~:
\[
  U \cap \bigcup_{i\in I} V_i = \bigcup_{i\in I}(U \cap V_i).
\]
\end{definition}

Ce qui rend cette construction cat\'egoriquement int\'eressante, c'est qu'elle
est fonctorielle~: les applications continues induisent des morphismes de
cadres, mais en sens inverse.

\begin{proposition}[Fonctorialit\'e de $\mathcal{O}$]
L'assignation $X \mapsto \mathcal{O}(X)$ s'\'etend en un foncteur contravariant
$\mathcal{O}\colon \mathbf{Top} \to \mathbf{Frm}^{\mathrm{op}}$,
o\`u $\mathbf{Frm}$ est la cat\'egorie des cadres. Pour $f\colon X \to Y$ continue,
$\mathcal{O}(f) = f^{-1}\colon \mathcal{O}(Y) \to \mathcal{O}(X)$ est un morphisme
de cadres.
\end{proposition}

\begin{proof}
La pr\'eimage $f^{-1}$ pr\'eserve les r\'eunions quelconques et les intersections finies,
donc c'est bien un morphisme de cadres. La fonctorialit\'e d\'ecoule de
$(g \circ f)^{-1} = f^{-1} \circ g^{-1}$.
\end{proof}

\begin{remarque}
Ce foncteur est au c\oe ur de la th\'eorie des locales (Chapitre~3), o\`u l'on renverse
la perspective~: au lieu d'\'etudier un espace via ses points, on l'\'etudie via son treillis
d'ouverts.
\end{remarque}

% ------------------------------------------------------------
\section{Topologie initiale et topologie finale}
% ------------------------------------------------------------

Pour construire de nouvelles topologies, deux strat\'egies s'opposent. La
topologie \emph{initiale} est la plus petite topologie qui rend certaines
applications continues~: c'est ainsi que l'on construit les produits et les
sous-espaces. La topologie \emph{finale} est la plus grande~: c'est le
m\'ecanisme derri\`ere les quotients et les sommes disjointes. Ces deux
constructions sont duales, et leurs propri\'et\'es universelles encapsulent
tout ce qu'il faut savoir pour les manipuler.

\begin{definition}[Topologie initiale]
Soit $X$ un ensemble et $(f_i\colon X \to (Y_i,\tau_i))_{i\in I}$ une famille
d'applications. La \emph{topologie initiale} sur $X$ est la topologie la moins fine
rendant toutes les $f_i$ continues. Elle est engendr\'ee par la sous-base
$\{f_i^{-1}(U) : i \in I, U \in \tau_i\}$.
\end{definition}

\begin{definition}[Topologie finale]
Soit $Y$ un ensemble et $(f_i\colon (X_i,\tau_i) \to Y)_{i\in I}$ une famille
d'applications. La \emph{topologie finale} sur $Y$ est la topologie la plus fine
rendant toutes les $f_i$ continues~:
$\tau_Y = \{V \subseteq Y : f_i^{-1}(V) \in \tau_i \text{ pour tout } i\}$.
\end{definition}

La puissance de ces constructions r\'eside dans leurs propri\'et\'es
universelles, qui r\'eduisent la v\'erification de la continuit\'e \`a un test
composante par composante.

\begin{theoreme}[Propri\'et\'e universelle]
\begin{enumerate}
  \item \emph{Topologie initiale~:} $g\colon Z \to X$ est continue si et seulement si
    $f_i \circ g$ est continue pour tout $i$.
  \item \emph{Topologie finale~:} $g\colon Y \to Z$ est continue si et seulement si
    $g \circ f_i$ est continue pour tout $i$.
\end{enumerate}
\end{theoreme}

\begin{proof}
(1) Si $g$ est continue, les $f_i \circ g$ le sont comme compos\'ees de fonctions continues.
R\'eciproquement, pour tout $U \in \tau_i$, $(f_i \circ g)^{-1}(U) = g^{-1}(f_i^{-1}(U))$
est ouvert dans $Z$. Comme les $f_i^{-1}(U)$ forment une sous-base de la topologie
initiale, $g$ est continue.

(2) Si $g$ est continue, les $g \circ f_i$ le sont. R\'eciproquement, soit $V$ ouvert dans
$Z$. Alors pour tout $i$, $f_i^{-1}(g^{-1}(V)) = (g \circ f_i)^{-1}(V)$ est ouvert dans
$X_i$, donc $g^{-1}(V)$ est ouvert dans la topologie finale.
\end{proof}

% ------------------------------------------------------------
\section{Espaces de Kolmogorov et quotient $T_0$}
% ------------------------------------------------------------

Revenons aux axiomes de s\'eparation pour une construction qui illustre
magnifiquement la pens\'ee cat\'egorielle. \'Etant donn\'e un espace topologique
quelconque, peut-on le \og rendre\fg{} $T_0$ de la mani\`ere la plus
\'economique possible~? La r\'eponse est le \emph{quotient de Kolmogorov}, et
c'est un exemple d'\emph{adjonction} en action.

\begin{definition}[Quotient de Kolmogorov]
Soit $(X,\tau)$ un espace topologique. On d\'efinit la relation d'\'equivalence~:
$x \sim y$ si et seulement si $\overline{\{x\}} = \overline{\{y\}}$. Le \emph{quotient
de Kolmogorov} est $X_0 = X / {\sim}$ muni de la topologie quotient.
\end{definition}

L'id\'ee est simple~: on identifie les points que la topologie ne peut pas
distinguer. Le r\'esultat est un espace $T_0$, et cette construction est
\og la meilleure possible\fg{} au sens pr\'ecis suivant.

\begin{theoreme}[R\'eflexion $T_0$]
Le quotient de Kolmogorov $X_0$ est un espace $T_0$, et l'application quotient
$q\colon X \to X_0$ est universelle~: pour toute application continue $f\colon X \to Y$
avec $Y$ un espace $T_0$, il existe une unique application continue
$\tilde{f}\colon X_0 \to Y$ telle que $f = \tilde{f} \circ q$.
Cat\'egoriquement, le foncteur d'inclusion $\mathbf{Top}_0 \hookrightarrow \mathbf{Top}$
poss\`ede un adjoint \`a gauche.
\end{theoreme}

\begin{proof}
Montrons que $X_0$ est $T_0$. Soient $[x] \neq [y]$ dans $X_0$. Alors
$\overline{\{x\}} \neq \overline{\{y\}}$, donc il existe un ouvert $U$ de $X$
contenant l'un mais pas l'autre. Comme $U$ est satur\'e pour $\sim$ (car si
$z \in U$ et $\overline{\{z\}} = \overline{\{z'\}}$, alors $z' \in U$), $q(U)$
est ouvert dans $X_0$ et s\'epare $[x]$ de $[y]$.

L'universalit\'e~: si $f\colon X \to Y$ est continue et $Y$ est $T_0$, alors
$\overline{\{x\}} = \overline{\{y\}}$ implique $\overline{\{f(x)\}} = \overline{\{f(y)\}}$
(car $f$ est continue), donc $f(x) = f(y)$ car $Y$ est $T_0$. Ainsi $f$ passe au quotient.
\end{proof}

% ------------------------------------------------------------
\section{Compacit\'e et variantes}
% ------------------------------------------------------------

Terminons ce chapitre de rappels par une mise au point terminologique
cruciale. La notion de compacit\'e, qui est si centrale en topologie, fait
l'objet d'un d\'esaccord de convention entre les traditions fran\c{c}aise et
anglo-saxonne qui peut engendrer de s\'erieuses confusions.

\begin{definition}[Compacit\'e]
Un espace $X$ est \emph{compact} si de tout recouvrement ouvert de $X$, on peut extraire
un sous-recouvrement fini.
\end{definition}

\begin{attention}
Contrairement \`a la convention de Bourbaki, nous ne supposons \emph{pas} qu'un espace
compact est s\'epar\'e. Cette convention, standard en th\'eorie des domaines et en
topologie alg\'ebrique, est essentielle pour ce cours. Lorsque nous voudrons un espace
compact et s\'epar\'e, nous dirons explicitement <<~compact de Hausdorff~>>.
\end{attention}

\begin{definition}[Quasi-compacit\'e]
Certains auteurs (suivant Bourbaki) appellent \emph{quasi-compact} ce que nous appelons
compact, et r\'eservent <<~compact~>> pour <<~compact de Hausdorff~>>. Dans ce cours,
\textbf{compact = quasi-compact au sens de Bourbaki}.
\end{definition}

% ------------------------------------------------------------
\section{Exercices}
% ------------------------------------------------------------

\begin{exercice}
Montrer que l'espace de Sierpi\'nski $\mathbb{S}$ est le classifieur des ouverts~: pour
tout espace $X$, il y a une bijection entre $\mathcal{O}(X)$ et l'ensemble des applications
continues $X \to \mathbb{S}$.
\end{exercice}

\begin{exercice}
Montrer que dans la cat\'egorie $\mathbf{Top}$, l'\'egalisateur de $f,g\colon X
\rightrightarrows Y$ est le sous-espace $\{x \in X : f(x) = g(x)\}$ muni de la topologie
induite.
\end{exercice}

\begin{exercice}
V\'erifier que le foncteur $\mathcal{O}\colon \mathbf{Top}^{\mathrm{op}} \to \mathbf{Frm}$
est fid\`ele si et seulement si on se restreint aux espaces $T_0$. Montrer un
contre-exemple dans le cas non $T_0$.
\end{exercice}

\begin{exercice}
Soit $X$ un ensemble ordonn\'e fini. Montrer que la topologie d'Alexandrov sur $X$ (dont
les ouverts sont les parties hautes $U$ telles que $x \in U$ et $x \leq y$ impliquent
$y \in U$) est la plus fine topologie dont l'ordre de sp\'ecialisation est l'ordre donn\'e.
\end{exercice}

\begin{exercice}
Montrer que le foncteur d'oubli $U\colon \mathbf{Top} \to \mathbf{Set}$ ne pr\'eserve ni
les coproduits ni les co\'egalisateurs. Donner des exemples explicites.
\end{exercice}

\begin{exercice}[Limites dans $\mathbf{Top}_0$]
Montrer que $\mathbf{Top}_0$ est une sous-cat\'egorie r\'eflexive de $\mathbf{Top}$ et en
d\'eduire qu'elle admet toutes les limites et colimites. Pr\'eciser~: les limites dans
$\mathbf{Top}_0$ co\"incident-elles avec celles de $\mathbf{Top}$~?
\end{exercice}
